\section{Barnes--Hut}
Dans ce cas, on peut proposer une distribution hiérarchique du quadtree entre les processus MPI. Comme le nombre de processus est une puissance de quatre, on peut choisir une profondeur \(L\) telle que \(4^L = P\), où \(P\) est le nombre de processus. Alors, les boîtes de ce niveau sont réparties entre les processus, de sorte que chaque processus reçoive une sous-boîte ainsi que tout le sous-arbre qui en dépend. Les boîtes des niveaux plus élevés, comme la racine et les nœuds communs, peuvent être dupliquées dans tous les processus, ce que l’énoncé permet effectivement.

À chaque pas de temps, tous les processus calculent d’abord ensemble la boîte globale qui contient toutes les étoiles. Ensuite, chaque processus construit localement le sous-arbre de sa zone en utilisant les étoiles qui lui correspondent. Puis, une phase de communication avec MPI est réalisée afin d’échanger les informations importantes des nœuds de l’arbre, en particulier la masse totale et le centre de masse de chaque boîte. Ainsi, chaque processus peut disposer d’une vision globale du quadtree, du moins sous une forme résumée.

Après cela, le calcul des accélérations peut être parallélisé en répartissant les étoiles entre les processus selon cette même distribution spatiale. Chaque processus calcule l’accélération de ses propres étoiles. Lorsqu’une boîte lointaine satisfait le critère de Barnes--Hut, sa contribution est approximée en utilisant uniquement sa masse totale et son centre de masse. Si ce critère n’est pas satisfait, il faut alors descendre récursivement vers les sous-boîtes. Si l’on atteint une feuille contenant peu d’étoiles, on ne réalise plus d’approximation, mais on utilise directement les positions et les masses de ces étoiles pour effectuer le calcul.

Cette stratégie permet de paralléliser assez efficacement le calcul de l’accélération, mais elle entraîne également un coût supplémentaire de communication, puisqu’il est nécessaire de diffuser les informations de l’arbre entre les processus. De plus, un déséquilibre de charge peut apparaître si certaines sous-boîtes contiennent beaucoup plus d’étoiles que d’autres, par exemple celles qui sont les plus proches du centre de la galaxie.
